\documentclass[10pt,a4paper]{amsart}
\usepackage[margin=19mm]{geometry}
\usepackage{amsmath,amssymb,mathtools}
\usepackage[scr=rsfs,cal=euler]{mathalfa}
\usepackage{xfrac}
\usepackage[unicode,hidelinks,bookmarks=false]{hyperref}
\newcommand{\ie}{i.e.\ }
\newcommand{\cf}{cf.\ }
\newcommand{\resp}{respectively\ }
\newcommand{\Subsection}{\subsection}
\providecommand{\nobreakdash}{\nobreak\hbox{-}}
\setlength{\emergencystretch}{2em}
\multlinegap0pt
\usepackage{bm}
\usepackage{bbm}
\usepackage{dsfont}
\usepackage{xspace}
\usepackage{cancel}
\usepackage{mathtools}
\usepackage{amsthm}
\makeatletter
\@ifundefined{theorem}{\newtheorem{theorem}{Theorem}}{}
\@ifundefined{prop}{\newtheorem{prop}[theorem]{Proposition}}{}
\makeatother
\title[Asymmetric phase transitions in random noncommutative geomet]{Asymmetric phase transitions in random noncommutative geometries}
\author{Benedek Bukor, Masoud Khalkhali, Samuel Kov\'a\v{c}ik, Adam Kubi\'s, Katar\'ina Magdolenov\'a, Nathan Pagliaroli,, Juraj Tekel}
\date{Source: arXiv:2609.19898v1 (CC BY 4.0)}
\begin{document}
\maketitle

\noindent\emph{Source and licence.} Asymmetric phase transitions in random noncommutative geometries. Version arXiv:2609.19898v1 (CC BY 4.0), distributed under the Creative Commons Attribution 4.0 International licence, as recorded in the arXiv licence metadata for this exact version and shown on the abstract page https://arxiv.org/abs/2609.19898v1. Full rights notice: licenses/arxiv-2609.19898-CC-BY-4.0.txt.\par\smallskip

\noindent\textit{Editorial scope.} Counted unit: the Prop whose source label is \texttt{free energy theorem} in the article (source0.tex lines 1089--1125, source label \texttt{free energy theorem}), delivered together with the article's own complete original proof of it (source0.tex lines 1127--1211, one \texttt{proof} environment of 85 lines). No mathematics is written, repaired or re-derived: statement and proof are the article's own text line for line. Required context retained uncounted: thm: one-cut density (lines 194--201); eq:residue formula for moments one-cut (lines 765--769); eq:m1 one-cut (lines 781--783); eq:m2 one-cut (lines 785--788); eq:m3 one-cut (lines 789--793); eq:m4 one-cut (lines 794--799); eq:1-cut condition 2 (lines 802--812); eq:1-cut condition 1 (lines 814--825). Every definition, display and label of those lines is retained unchanged. Omitted from this package and not redistributed: 1--193, 202--764, 770--780, 784--784, 800--801, 813--813, 826--1088, 1126--1126, 1212--1218. Nothing inside a retained excerpt is omitted or altered: every retained body is delivered exactly as the article's lines are written, so no omission receipt of any kind accompanies this package. Citations: the retained excerpts carry no citation command at all, so no bibliography entry of the article is transcribed and no reference is redistributed. Reference adaptation: none was needed and none was made; every reference inside the delivered unit resolves inside it, so no unresolved reference or citation is produced. Printed slips kept literal: no printed word, symbol, hypothesis or number of the counted statement or of its proof was altered. Layout and class: the article's own class is \texttt{article} with packages including authblk, amsfonts, float, hyperref, caption, multirow, subcaption, lineno, amssymb, latexsym, amsmath, amscd, slashed, mathtools; the wrapper is the lane's pinned \texttt{amsart} editorial wrapper at 10pt on A4, which restates the theorem-like environments the retained text needs and adds the article's own macro definitions that the retained text needs (none), with extra wrapper packages (bm, bbm, dsfont, xspace, cancel, mathtools). The delivered numbering is the unit's own. Assets: no retained span contains a figure environment, a graphics-inclusion command, a PDF-page inclusion, a table or a TikZ picture; no image file is copied into this package, the package declares no asset dependency and no image file is redistributed with it. Licence: version arXiv:2609.19898v1 of this article is distributed under the Creative Commons Attribution 4.0 International licence, as recorded in the arXiv licence metadata for this exact version and shown on the abstract page https://arxiv.org/abs/2609.19898v1. A full rights notice is delivered at licenses/arxiv-2609.19898-CC-BY-4.0.txt. Characters: every retained excerpt is pure ASCII, so the delivered engine, XeLaTeX with its default Latin Modern text font, reports no missing character.
	\begin{theorem}\label{thm: one-cut density}
			The spectral density function of the one-cut solution of the type $(1,0)$  quartic Dirac ensemble is
		\begin{equation*}
		\rho(x) = \frac{1}{2\pi} \left(\frac{1}{\gamma ^2}-\frac{24 \epsilon t_{4} (x-\alpha ) \left(\alpha +24 \alpha \gamma ^4 t_4\right)}{24 \gamma ^4
 \epsilon t_{4}-1}+8 t_{4} \left(-2 \alpha ^2-\gamma ^2+x^2+\alpha x\right)\right)\sqrt{(x-a)(b-x)}_{[a,b]},
		\end{equation*}
 where $\alpha = \frac{a+b}{2}$ and $\gamma =\frac{b-a}{4}$ are the roots of the polynomials \eqref{eq:1-cut condition 2} and \eqref{eq:1-cut condition 1} and $\sqrt{\ldots}_\mathcal I$ denotes a function vanishing outside of the set of intervals $\mathcal I$.
	\end{theorem}

	\begin{align}\label{eq:residue formula for moments one-cut}
		m_{\ell} &= - \text{Res}_{x\rightarrow \infty} x^{\ell} W(x) dx\\
		&= - \text{Res}_{z\rightarrow \infty} \left(\alpha + \gamma\left(z+\frac{1}{z}\right) \right)^{\ell} W(x(z)) \left( \gamma\left(1-\frac{1}{z^2}\right)\right)dz\\
		&= \sum_{0\leq i,j\leq\ell} \frac{\ell!}{j! i!(\ell-i-j)!}\alpha^{\ell-i-j}\gamma^{i+j+1}(u_{i-j+1}-u_{i-j-1}).
	\end{align}

	\begin{equation}\label{eq:m1 one-cut}
		m_1 = \alpha \left(\frac{1+24 t_{4}\gamma^4}{1-24 t_{4}\epsilon\gamma^4}\right),
	\end{equation}

	\begin{align}\label{eq:m2 one-cut}
		m_{2} &=\gamma ^2-\frac{\alpha ^2 \left(24 \gamma ^4 (\epsilon+2) t_{4}+1\right)}{24 \gamma ^4
 \epsilon t_{4}-1}+8 \gamma ^6 t_{4}
 \end{align}

 \begin{align}\label{eq:m3 one-cut}
		m_3 &=\frac{3 \alpha \gamma ^2 \left(24 \gamma ^4 t_{4}+1\right) \left(8 \gamma ^4 \epsilon
 t_{4}-1\right)-\alpha ^3 \left(24 \gamma ^4 (2 \epsilon+3) t_{4}+1\right)}{24 \gamma ^4
 \epsilon t_{4}-1}
 \end{align}

 \begin{align}\label{eq:m4 one-cut}
 m_{4} &=\frac{6 \alpha ^2 \gamma ^2 \left(192 \gamma ^8 \epsilon t_{4}^2-8 \gamma ^4 (\epsilon+5)
 t_{4}-1\right)-\left(\alpha ^4 \left(24 \gamma ^4 (3 \epsilon+4) t_{4}+1\right)\right)+2
 \gamma ^4 \left(12 \gamma ^4 t_{4}+1\right) \left(24 \gamma ^4 \epsilon
 t_{4}-1\right)}{24 \gamma ^4 \epsilon t_{4}-1}.
	\end{align}

 \begin{align}\label{eq:1-cut condition 2}
 \begin{split}
		0&=\frac{4 \alpha}{1-24 \gamma ^4 \epsilon t_{4}} \left((\epsilon+1) t_{2}+2 t_{4} \left(3 \gamma ^2 \left(-2 \epsilon
 \left(96 \gamma ^8 t_{4}^2+12 \gamma ^4 t_{4}-1\right)+8 \gamma ^4
 t_{4}+3\right) \right.\right.\\
 &\left.\left.+\alpha ^3 \epsilon \left(24 \gamma ^4 (2 \epsilon+3)
 t_{4}+1\right)+\alpha ^2 \left(3 \epsilon \left(40 \gamma ^4 t_{4}+1\right)+144 \gamma
 ^4 t_{4}+4\right)-3 \alpha \gamma ^2 \epsilon \left(24 \gamma ^4 t_{4}+1\right) \left(8
 \gamma ^4 \epsilon t_{4}-1\right)\right)\right)
 \end{split}
	\end{align}

	\begin{align}\label{eq:1-cut condition 1}
 \begin{split}
		1&= \frac{4 \gamma^2}{24 \gamma ^4 \epsilon
 t_{4}-1} \left(2 t_{4} \left(-\left(\alpha ^3 \left(24 \gamma ^4 \epsilon (2
 \epsilon+3) t_{4}+\epsilon\right)\right)-6 \alpha ^2 (\epsilon+1) \left(24 \gamma ^4
 t_{4}+1\right)\right. \right.\\
 &\left. \left.+3 \alpha \gamma ^2 \epsilon \left(24 \gamma ^4 t_{4}+1\right) \left(8
 \gamma ^4 \epsilon t_{4}-1\right)+6 \gamma ^2 \left(4 \gamma ^4 t_{4}+1\right) \left(24
 \gamma ^4 \epsilon t_{4}-1\right)\right)+t_{2}(24 \gamma ^4 \epsilon
 t_{4}-1)\right),
 \end{split}
	\end{align}

\begin{prop}\label{free energy theorem}
 	The free energy for a one-cut solution of the type $(1,0)$ quartic Dirac ensemble, assuming $0<a<b$ is of the form 
 \begin{equation}
 I[\rho_{\textit{1-cut}}] = E(t_{2},t_{4}) +2t_{4}(m_{4} + 3 m_{2}^2 + 4 \epsilon m_{1}m_{3}) + 2t_{2}(m_{2}+ \epsilon m_{1}^2),
 \end{equation}
 where
 \begin{itemize}
 \item the first term is \begin{align*}&E(t_{2},t_{4})=
\gamma^2\left( 8 t_{4}\alpha^2+B\alpha+C\right)
\left(
\ln\left(\frac{q}{2}\right)+\frac{2\gamma^2}{q^2}
\right)
+
\frac{2\gamma^4(16 t_{4}\alpha+B)}{q}
\left(
1-\frac{4\gamma^2}{3q^2}
\right)
+
8 t_{4}\gamma^4
\left(
\ln\left(\frac{q}{2}\right)+\frac{4\gamma^4}{q^4}
\right),
 \end{align*}
 where $q = \alpha + \sqrt{\alpha^2 - 4 \gamma^2}$,
 \begin{align*}
 B& = 8 \alpha t_{4}-\frac{24 \epsilon t_{4} \left(\alpha +24 \alpha \gamma ^4 t_{4}\right)}{24 \gamma ^4 \epsilon
 t_{4}-1},
 \end{align*}
 and
 \begin{align*}
 C & = \frac{1}{\gamma ^2}+\frac{24 \alpha \epsilon t_{4} \left(\alpha +24 \alpha \gamma ^4 t_{4}\right)}{24 \gamma ^4
 \epsilon t_{4}-1}-16 \alpha ^2 t_{4}-8 \gamma ^2 t_{4}.
 \end{align*}
 \item $m_{1},m_{2},$ $m_{3},$ and $m_{4}$ are as in (\ref{eq:m1 one-cut}-\ref{eq:m4 one-cut}) 
 \item $\gamma = \frac{b-a}{4}$ and $\alpha=\frac{a+b}{2}$, and are the solutions of equations \eqref{eq:1-cut condition 1} and \eqref{eq:1-cut condition 2}.
 \end{itemize}
\end{prop}

\begin{proof}
As in the previous proof, we wish to compute integrals of the form 
\begin{equation*}
 L_{\ell}:= \int_{a}^{b} \ln|x| x^\ell \sqrt{(x-a)(b-x)}dx.
\end{equation*}
for $\ell = 0,1,$ and 2. 


Applying the substitution $a=s-t$ and $b= s+t$, we can write the discriminant as $t^2 - (x-s)^2$. Then making the change of variables $x \rightarrow s+t\cos (\theta)$ we have that
\begin{align*}
 L_{\ell}
 &=\frac{t^2}{2}\int_{0}^{\pi} \ln|s+t\cos (\theta)| (s+t\cos (\theta))^\ell \left(1-\cos(2 \theta)\right) d \theta.
\end{align*}
We may write the integrand as a Fourier cosine series 
\begin{equation*}
 \ln|s+ t \cos(\theta)| =  \sum_{n=0}^{\infty} a_{n}\cos( n \theta).
\end{equation*}



To compute the Fourier coefficients we start by mapping the interval $[0,\pi]$ to the unit circle with $z = e^{i\theta}$. Let $\alpha = \frac{-s+\sqrt{s^2-t^2}}{t}$ and $\beta = \frac{-s-\sqrt{s^2-t^2}}{t}= \frac{1}{\alpha}$ denote the roots of the argument of the logarithm after this transformation. 

For the case that $\ell=0$, we have that 
$$
 \frac{t^2}{2}\int_0^\pi\left(\frac{a_0}{2}+\sum_{n \geq 1} a_n \cos (n \theta)\right)\left(1-\cos (2 \theta)\right) d \theta= \frac{\pi t^2}{4}\left(a_0-a_2\right).
$$
Then we may use the Residue Theorem to compute
\begin{align*}
 a_{0} &= \frac{1}{\pi}\int_{0}^{\pi}\ln|s+t \cos(\theta)|d\theta \\
 &=\frac{1}{\pi i} \oint_{|z|=1}\ln\left| s+ \frac{t}{2}\left(z+\frac{1}{z} \right)\right|\frac{dz}{z}\\
 &=2\text{Res}_{z\rightarrow 0}\left[-\frac{1}{z^2}\ln\left( s+ \frac{t}{2}\left(z+\frac{1}{z} \right)\right)\right]\\
 &= 2 \text{Res}_{z\rightarrow 0}\left[\frac{1}{z}\ln \left(\frac{s+\sqrt{s^2-t^2}}{2}\right)-\frac{1}{z}\left(\sum_{n=1}^{\infty} \frac{\alpha^n}{n z^n}+\sum_{n=1}^{\infty} \frac{\alpha^n z^n}{n} \right)\right]\\
 &= 2\ln \left(\frac{s+\sqrt{s^2-t^2}}{2}\right),
\end{align*}
and via the same approach
\begin{align*}
 a_{n} = \frac{2(-1)^{n+1}}{n} r^{n},
\end{align*}
where $r = \frac{t}{s + \sqrt{s^2 -t^2}}$, for $n\geq 1$. Thus
\begin{equation*}
 L_{0} = \frac{\pi t^2}{4} \left( 2\ln \left(\frac{s+\sqrt{s^2-t^2}}{2}\right) +r^2 \right).
\end{equation*}


When $\ell =1$, we have 
\begin{align*}
 L_{1} &=\frac{st^2}{2}\int_{0}^{\pi} \ln|s+t\cos (\theta)| \left(1-\cos(2 \theta)\right) d \theta
 + \frac{t^3}{2}\int_{0}^{\pi} \ln|s+t\cos (\theta)| \cos (\theta)\left(1-\cos(2 \theta)\right) d \theta.
\end{align*}
The Fourier cosine series expansion of the second term is 
$$
 \frac{t^2}{2}\int_0^\pi\left(\frac{a_0}{2}+\sum_{n \geq 1} a_n \cos ( n \theta)\right)\cos(\theta)\left(1-\cos (2\theta)\right) d \theta= \frac{\pi t^2 }{8}(a_{1} -a_{3}).
$$
Hence, 
\begin{equation*}
 L_{1} = s L_{0} + \frac{\pi t^{3}}{4}\left(r - \frac{r^3}{3} \right).
\end{equation*}

When $\ell =2$, we have 
\begin{align*}
 L_{2} &=\frac{t^2s^2}{2}\int_{0}^{\pi} \ln|s+t\cos (\theta)| \left(1-\cos(2 \theta)\right) d \theta +st^3\int_{0}^{\pi} \ln|s+t\cos (\theta)|\cos(\theta) \left(1-\cos(2 \theta)\right) d \theta\\
 &+\frac{t^4}{2}\int_{0}^{\pi} \ln|s+t\cos (\theta)| \cos(\theta)^2\left(1-\cos(2 \theta)\right) d \theta.\\
\end{align*}
The Fourier cosine series for the third integral is 
\begin{equation*}
 \frac{t^{4}}{2}\int_0^\pi\left(\frac{a_0}{2}+\sum_{n \geq 1} a_n \cos (n \theta)\right)\cos(\theta)^2\left(1-\cos (2 \theta)\right) d \theta= \frac{\pi t^{4}}{16}(a_{0}-a_{4}).
\end{equation*}
Hence, 
\begin{equation*}
 L_{2} = s^2 L_{0} + \frac{\pi s t^3}{2}\left( r - \frac{r^3}{3}\right) + \frac{\pi t^{4}}{8}\left(\ln \left(\frac{s+\sqrt{s^2-t^2}}{2}\right) - \frac{r^{4}}{4}\right).
\end{equation*}



For a density 
$\rho(x) = \frac{1}{2\pi}(Ax^2 + Bx + C)\sqrt{(x-a)(b-x)}$, we can compute 
\begin{align*}
 \int_{a}^{b} \ln|x| \rho(x) dx &= \frac{1}{2 \pi}( AL_{2} + B L_{1} + CL_{0}).
\end{align*}
Applying the inverse change of variables $s =\frac{a+b}{2}$ and $t =\frac{b-a}{2}$, and using the appropriate choice of $A$, $B$, and $C$ from Theorem \ref{thm: one-cut density} gives us the formula for $E(t_{2},t_{4})$.



Lastly, to compute the integral of $U(x,y)\rho(x)\rho(y)$ we can use equations \eqref{eq:m1 one-cut},\eqref{eq:m2 one-cut},\eqref{eq:m3 one-cut} for the moments $m_{1},m_{2}$ and $m_{3}$, but we also require $m_{4}$, which can be computed in the same manner using the residue formula \eqref{eq:residue formula for moments one-cut}.
\end{proof}

\end{document}
